% r04 Walled tori, the closing estimate and rotation invariance
\section{Walled tori, the closing estimate and rotation invariance}

\subsection{Walled tori and the closing estimate}

\begin{definition}[admissible walls and walled values]\label{rot:def:walls}
For a bounded set $X\subset\R^3$ let $\Aw{X}$ be the set of walled configurations $\fc=(J',w)$ of \cref{rot:def:tori} with
$\Nr{1}{X}\subset(-25MJ',25MJ')^3$ and $w=0$ on $\Nr{1}{X}$. The interior $(-25MJ',25MJ')^3$ of the cell $\Gamma_{J'}$ is a
fundamental domain of $D_{J'}$ for both lattices, $R^{-1}\Lk$ descending to $D_{J'}$ for $R\in\Rgen$ by \cref{rot:lem:geom}(b), and
sources supported in $X$ are periodised from it. If $X\subset Y$ then $\Aw{Y}\subset\Aw{X}$. Recall the walled ratio
$\Tw{a}{\fc,\Lambda}(f,g)=e^{-a_k\sS_\Lambda(g)}\langle\Theta(\mathrm if,g)\rangle_{\fc,\Lambda}$ of a lattice $\Lambda$ of spacing
$\eps_k$ descending to $D_{J'}$, defined for $(f,g)\in\cK$ with $w=0$ on $\supp g$ (\cref{rot:lem:walledalg}(c)). For $g\ne0$ assume \Ka\ (\Cref{rot:hyp:Ka}). For
$(f,g)\in\cK$ supported in $X$ and $\fc\in\Aw{X}$ the periodised sources are smooth, and by \cref{rot:lem:K}(a) the \emph{walled value}
\[
\Wv^a_\fc(f,g)\defeq\lim_{k\to\infty}\Tw{a}{\fc,\Lk}(f,g)=\lim_{k\to\infty}\Tw{a}{\fc,R^{-1}\Lk}(f,g)\qquad(R\in\Rgen)
\]
exists in $(0,\infty)$. For $b>0$ and $(f,g)\in\cK$ supported in $\Qb{b}$ let $I^a_b(f,g)\defeq\inf_{\fc\in\Aw{\Qb b}}\Wv^a_\fc(f,g)$,
which is nondecreasing in $b$, and
\[
\Wenva(f,g)\defeq\lim_{b\to\infty}I^a_b(f,g)\in(0,\infty].
\]
At $g=0$ these definitions rest on \Kz\ alone (\Cref{rot:hyp:K}; \cref{rot:lem:K}(a) at the complex coefficient $\mathrm i$), and they do not depend on the
scheme (\cref{rot:lem:walledalg}(c)): $\Wv_\fc(f,0)=\lim_k\langle e^{\mathrm iF_{\Lk}(f)}\rangle_{\fc,\Lk}\in[-1,1]$, and we write
$I_b(f,0)$, $\Wenv(f,0)$. King's printed state at a curvature source is written $\ZKinf(\mathrm if,0)$ in this part, the second slot
recording the absence of a scalar source; it is $\ZKinf(\mathrm if)$ of \cref{thm:s05:CK}(b), and $\ZKJ(\mathrm if)$ denotes the fixed-cube limit there.
\end{definition}

\begin{lemma}[the free direction]\label{rot:lem:free}
\begin{enumerate}[label=\textup{(\alph*)},leftmargin=*]
\item Assume $\Wadm{a}$ and \Ka. If $(f,g)\in\cK$ is supported at $\ell^\infty$-distance $>1$ from $\R^3\setminus\CJ$ and
$\fc\in\Aw{\bar C_J}$, then $\Wv^a_\fc(f,g)\ge Z^a_J(\mathrm if,g)$, where $Z^a_J=\lim_kZ^a_{k,J}$ is the fixed-cube limit of
\cref{rot:thm:full}(i).
\item Assume $\Wadm{a}$ and \Ka. Then $\Wenva\ge Z^a_\infty$ on $\cK$.
\item If the real $2$-form $f\in C^\infty_c$ is supported at $\ell^\infty$-distance $>1$ from $\R^3\setminus\CJ$ and
$\fc\in\Aw{\bar C_J}$, then $\Wv_\fc(f,0)\ge\ZKJ(\mathrm if)$.
\item $\Wenv(f,0)\ge\ZKinf(\mathrm if,0)\ge0$ for every real $2$-form $f\in C^\infty_c$.
\end{enumerate}
\end{lemma}

\begin{proof}
(a) Fix $k$. The set $\XJ=\CJ\cap\Lk$ and its nearest neighbours lie in $\Nr{1}{\bar C_J}$, hence in the interior of the cell, and the wall
vanishes there; so \cref{rot:lem:cutwall} applied to the walled torus state $\langle\cdot\rangle_{\fc,\Lk}$ with $X=\XJ$ has as its
$X$-factor the comparison cube $\ava{\cdot}$ of \cref{def:s02:acube}, and gives $\ava{\Theta(\mathrm if,g)}\le\langle\Theta(\mathrm if,g)\rangle_{\fc,\Lk}$;
every plaquette and site at which $(f,g)$ does not vanish lies in $\XJ$ by the support condition. The multiplier $e^{-a_k\sS_{\Lk}(g)}$
is the same positive constant on both sides (\cref{rot:lem:walledalg}(c), \cref{lem:s02:alg}(a)), so
$Z^a_{k,J}(\mathrm if,g)\le\Tw{a}{\fc,\Lk}(f,g)$. Let $k\to\infty$: the left side converges to $Z^a_J(\mathrm if,g)$ by
\cref{rot:thm:full}(i), the right side to $\Wv^a_\fc(f,g)$ by \cref{rot:def:walls}.

(b) If $\bar C_J\subset\Qb{b}$ then $\Aw{\Qb b}\subset\Aw{\bar C_J}$, so $I^a_b(f,g)\ge Z^a_J(\mathrm if,g)$ by (a) for every $J$ at which the
support condition of (a) holds; let $b\to\infty$ and then $J\to\infty$ (\cref{rot:thm:full}(iii)).

(c) At $g=0$ no multiplier enters. As in (a), \cref{rot:lem:cutwall} gives $Z_{k,J}(\mathrm if)=\ava{e^{\mathrm iF(f)}}\le T_{\fc,\Lk}(f,0)$,
and the first inequality of \cref{lem:s04:G}, whose support condition \eqref{eq:s04:S} holds for $\XJ$ because $f$ is supported at
distance $>1$ from $\R^3\setminus\CJ$, gives $Z^K_{k,J}(\mathrm if)=\avK{e^{\mathrm iF(f)}}\le Z_{k,J}(\mathrm if)$. Hence
$Z^K_{k,J}(\mathrm if)\le T_{\fc,\Lk}(f,0)$ at every $k$. Let $k\to\infty$: the left side converges to $\ZKJ(\mathrm if)$ by
\cref{thm:s05:CK}(b) under \EC, the right side to $\Wv_\fc(f,0)$ by \Kz\ (\cref{rot:lem:K}(a)).

(d) Fix $f$ and $J\ge J_0(f)$. For $b\ge MJ$ we have $\bar C_J=\Qb{MJ}\subset\Qb b$, so $\Aw{\Qb b}\subset\Aw{\bar C_J}$ and
$I_b(f,0)\ge\ZKJ(\mathrm if)$ by (c); letting $b\to\infty$, $\Wenv(f,0)\ge\ZKJ(\mathrm if)$ for every $J\ge J_0(f)$. By
\cref{cor:s04:G1} and \cref{thm:s05:CK}(b), $J\mapsto\ZKJ(\mathrm if)$ is nondecreasing with values in $[0,1]$, and by
\cref{cor:s04:G2}(ii)($\beta$) it converges to $\ZKinf(\mathrm if,0)$; let $J\to\infty$.
\end{proof}

\Cref{rot:lem:free}(a),(c) compare a walled value with a cube of fixed size. They do not assert $\Wv^a_\fc\ge Z^a_\infty$ for an
individual configuration, and this need not hold: for instance, at $e=\lambda=0$, $f=0$ and $g\ge0$ small, as $n\to\infty$ the walled
value of $(J',\psi,n)$ decreases (\cref{lem:s02:BFS}(c)) towards the value of the Dirichlet state on the set where $\psi=0$, which is
smaller than the infinite-volume value, as the Gaussian formula for $\langle e^{\Wick{g}}\rangle$ shows.

\begin{lemma}[independence of the orientation]\label{rot:lem:Rind}
Assume \Ka\ if $g\ne0$. For $R\in\Rgen$ and $(f,g)\in\cK$,
\[
\lim_{L\to\infty}\ \inf_{\fc\in\Aw{R^{-1}\Qb{ML}}}\Wv^a_\fc(f,g)=\Wenva(f,g),
\]
the infimum being defined as soon as $\supp f\cup\supp g\subset\Qb{5ML/7}$. At $g=0$ the same holds, with no assumption, for $\Wv_\fc(f,0)$ and
$\Wenv(f,0)$.
\end{lemma}

\begin{proof}
By \cref{rot:lem:geom}(c), $\Qb{5ML/7}\subset R^{-1}\Qb{ML}\subset\Qb{7ML/5}$, hence
$\Aw{\Qb{7ML/5}}\subset\Aw{R^{-1}\Qb{ML}}\subset\Aw{\Qb{5ML/7}}$, and
$I^a_{5ML/7}\le\inf_{\Aw{R^{-1}\Qb{ML}}}\Wv^a_\fc\le I^a_{7ML/5}$. Both bounds tend to $\Wenva$ as $L\to\infty$.
\end{proof}

\begin{theorem}[\TIz, \TItwo: the closing estimate]\label{rot:hyp:TIGHT}
\textbf{\TIz}: $\Wenv(f,0)\le\ZKinf(\mathrm if,0)$ for every real $2$-form $f\in C^\infty_c$.
\textbf{\TItwo}: for every real $2$-form $f\in C^\infty_c$,
\[
\lim_{b\to\infty}\ \sup_{\fc\in\Aw{\Qb b}}\sigma^2_\fc(f)\ \ge\ S^K_2(f),\qquad
\sigma^2_\fc(f)\defeq\lim_k\langle F_{\Lk}(f)^2\rangle_{\fc,\Lk},\quad
S^K_2(f)\defeq S^K_2(f,f)=\lim_J\lim_k\avK{F(f)^2}.
\]
\end{theorem}

\TIz\ involves only \Kz, which defines its left side, and King's printed state, which exists by \EC. It is proved in the heavy-shell sections: \Cref{rot:thm:REStoTIGHT} derives it from the decoupling property \REz\ of \Cref{rot:hyp:RES}, which \Cref{rot:thm:shell} establishes at one configuration from the heavy-shell estimates \HS; \TItwo\ follows from \TIz\ by \Cref{rot:lem:rel}(c). In the Wick sector the closing statement is a property of the scheme, not proved here (it is listed among the open problems, in the final section) and used only by the optional extension of the Wick leg (\Cref{thm:s23:FormA}):

\begin{definition}[\TIa: the closing property, Wick sector]\label{rot:hyp:TIa}
Let $a$ be $\mathrm W$-admissible and assume \Ka. \textbf{\TIa}: $\Wenva\le Z^a_\infty$ on $\cK$. Equivalently: for every $b>0$, every $(f,g)\in\cK$ supported in $\Qb b$ and every $\eta>0$
there is $\fc\in\Aw{\Qb b}$ with $\Wv^a_\fc(f,g)\le Z^a_\infty(\mathrm if,g)+\eta$.
\end{definition}

The equivalence in \TIa\ holds because $I^a_b$ is nondecreasing in $b$. By \cref{rot:lem:free}, \TIa\ is equivalent to
$\Wenva=Z^a_\infty$ on $\cK$, and \TIz\ to $\Wenv(\cdot,0)=\ZKinf(\mathrm i\,\cdot,0)$: \emph{King's printed state is the
limit of the lower envelopes of walled tori, the continuum limit being taken first}. The limit defining $\sigma^2_\fc$ exists by \Kz,
since $t\mapsto\langle e^{\mathrm itF(f)}\rangle_{\fc,\Lk}$ converges locally uniformly as an entire function (\cref{rot:lem:K}(a)) and
second derivatives of locally uniformly convergent entire functions converge; $S^K_2$ is the two-point Schwinger function of
\cref{thm:s05:CK}(c), with the lattice second moments converging to it. The supremum over $\Aw{\Qb b}$ is nonincreasing in $b$, so the
limit in \TItwo\ exists.

\begin{remark}[the configuration class]\label{rot:rem:class}
The proofs below use \Kz\ only for the walled configurations that enter them, and they use the infimum over $\Aw{R^{-1}\Qb{ML}}$ only
through an upper bound by the walled value of one configuration supplied by \TIz. Hence \Kz\ and \TIz\ may be read on any common
sub-class of walled configurations which is stable under raising the wall height and contains the single-level configuration produced
by \cref{rot:thm:REStoTIGHT}: the infimum is then taken over that sub-class, \TIz\ is read with the same sub-class, and
\cref{rot:lem:Rind} holds on it, since its proof uses only the inclusions between shrinking families. The sub-class on which the proof of
\Kz\ along King's lines \cite{K2} supplies the comparison consists of tori whose side is a large multiple of a power of $5$, of
mass-increasing walls whose mean over the torus exceeds a fixed multiple of the torus volume, and of a source box and torus aligned with
the blocks of the final scale. The configuration of \cref{rot:thm:REStoTIGHT} lies in it: its wall is $(n+K)\psi$ with $n+K\ge1$,
mass-increasing; the condition on the wall mean is linear in the height and holds once $K$ is large at a fixed torus, which
\cref{rot:rem:Eproof}(iii) allows; the torus and the alignment of the source box are choices made in \cref{rot:thm:shell}.
\end{remark}

\begin{lemma}\label{rot:lem:rel}
\begin{enumerate}[label=\textup{(\alph*)},leftmargin=*]
\item If $a,a'$ are $\mathrm W$-admissible and \Ka\ holds, then \TIa\ and $\mathrm{(TIGHT)}_{a'}$ are equivalent.
\item Assume $\Wadm{a}$ and \EC. Then \TIa\ implies \TIz.
\item \TIz\ implies \TItwo; and $\lim_b\sup_{\Aw{\Qb b}}\sigma^2_\fc(f)\le S^K_2(f)$ always, so \TItwo\ is an
equality.
\end{enumerate}
\end{lemma}

\begin{proof}
(a) By \cref{rot:lem:schemes}(b), $a'_k-a_k\to\kappa_{a'a}$, and $\Tw{a'}{\fc,\Lambda}=e^{-(a'_k-a_k)\sS_\Lambda(g)}\Tw{a}{\fc,\Lambda}$
by \cref{rot:lem:walledalg}(c); since $\sS_{\Lk}(g)\to\int g$, $\Wv^{a'}_\fc=e^{-\kappa_{a'a}\int g}\,\Wv^a_\fc$ with a factor independent
of $\fc$. The same factor relates $Z^{a'}_\infty$ and $Z^a_\infty$ (\cref{rot:thm:full}(vi)).

(b) At $g=0$ the walled values are scheme-free (\cref{rot:lem:walledalg}(c)), so the left sides of \TIa\ and \TIz\ agree, and the right
side of \TIa\ is $Z^a_\infty(\mathrm if,0)=\lim_J\lim_kZ_{k,J}(\mathrm if)$, the iterated limit of the comparison family, which exists
under $\Wadm{a}$ (\cref{rot:thm:full}(iii)) and equals $\ZKinf(\mathrm if,0)$ under \EC\ by \cref{cor:s04:G2}(ii)($\gamma$).

(c) Fix $f$, $b$ and $\eta'>0$. All functions of $t$ below are real and even. Since $\cos x\ge1-\frac{x^2}2$ and
$\cos x\le1-\frac{x^2}2+\frac{x^4}{24}$ at every lattice level, and the moments converge in the limits (for the walled torus by
\cref{rot:lem:K}(a), for King's cube by \cref{thm:s05:CK}(c)),
\[
\Wv_\fc(tf,0)\ge1-\tfrac{t^2}2\sigma^2_\fc(f),\qquad
\ZKinf(\mathrm itf,0)\le1-\tfrac{t^2}2S^K_2(f)+\tfrac{t^4}{24}M_4(f),
\]
with $M_4(f)\defeq\partial_t^4\ZKinf(tf,0)|_{t=0}=\lim_J\lim_k\avK{F(f)^4}<\infty$. By \TIz\ applied to $tf$, for every $t>0$ and
$\eta>0$ there is $\fc\in\Aw{\Qb b}$ with $\Wv_\fc(tf,0)\le\ZKinf(\mathrm itf,0)+\eta$, hence
$\sigma^2_\fc(f)\ge S^K_2(f)-\frac{t^2}{12}M_4(f)-\frac{2\eta}{t^2}$. Choosing $t$ with $t^2M_4(f)/12<\eta'/2$ and then
$\eta=t^2\eta'/4$ gives $\sup_{\Aw{\Qb b}}\sigma^2_\fc(f)\ge S^K_2(f)-\eta'$.

For the second claim, let $\fc\in\Aw{\bar C_J}$ with $f$ supported at distance $>1$ from $\R^3\setminus\CJ$. By
\cref{rot:lem:free}(c) applied to $tf$, $\Wv_\fc(tf,0)\ge\ZKJ(\mathrm itf)$ for all real $t$, with equality ($=1$) at $t=0$; both sides
are even and smooth, so comparing second derivatives at $t=0$ gives $\sigma^2_\fc(f)\le v_J(f)\defeq\lim_k\avK{F(f)^2}$, the limit
existing by \cref{thm:s05:CK}(c). Hence $\sup_{\Aw{\Qb b}}\sigma^2_\fc(f)\le v_J(f)$ whenever $\bar C_J\subset\Qb b$. By
\cref{cor:s04:G1}, $\avK{e^{\mathrm itF(f)}}\le\langle e^{\mathrm itF(f)}\rangle^K_{k,J+1}$ for all real $t$, with equality at $t=0$;
comparing second derivatives, the cube variances $\avK{F(f)^2}$ are nonincreasing in $J$ at every $k$, hence so are the $v_J(f)$, and
they converge to $S^K_2(f)$ (\cref{thm:s05:CK}(c)). Therefore $\lim_b\sup_{\Aw{\Qb b}}\sigma^2_\fc(f)\le\inf_Jv_J(f)=S^K_2(f)$.
\end{proof}

\begin{remark}[the role of the closing estimate]\label{rot:rem:gauss}
Apart from the closing estimate, the proofs of \cref{rot:thm:gen,rot:thm:curv} below use only: monotonicity of the cube functionals in
$J$ (on $\cK$ and on the decreasing class), the comparison ``cube $\le$ walled torus'' on $\cK$ and ``cube $\ge$ walled torus'' on the
decreasing class on each lattice, the equality of walled values on the two lattices, passive covariance, and the existence of the
limits. All of these hold for the following family of Gaussian functionals, in which the analogue of King's state is not rotation
invariant. Let $A=\diag(1,2,3)$, $\mu>0$, and let $C_I$, $C_R$, $C_U$ be the translation-invariant operators on $L^2(\R^3)$ with
Fourier symbols $\hat c_I(k)=(\ip k{Ak}+\mu^2)^{-1}$, $\hat c_R(k)=\hat c_I(Rk)$ and $\hat c_U=\min(\hat c_I,\hat c_R)$. Let the
``cube of size $J$'' on the lattice $X\in\{I,R\}$ have covariance $C_X+J^{-1}$, let every ``walled torus'' have covariance $C_U$ on
both lattices, and let a covariance $C$ define the functionals $e^{-\frac12\ip f{Cf}}$ on the cone and $e^{\frac12\ip f{Cf}}$ on the
decreasing class. Every relation listed reduces to $C_U\le C_X+J^{-1}$ or to the monotonicity of $C_X+J^{-1}$ in $J$, and passive
covariance holds since $\ip f{C_Rf}=\ip{R\cdot f}{C_I\,R\cdot f}$. But $e^{-\frac12\ip{R\cdot f}{C_IR\cdot f}}\ne e^{-\frac12\ip f{C_If}}$
in general, and the closing estimate fails since $C_U\ne C_I$. Hence the closing estimate is not a consequence of the other
properties, and an input asserting ``walled torus $\le$ cube $+$ small'' is needed. (The same family shows that any placement of the
two-lattice comparison on a finite box must also carry an input beyond the correlation inequalities: there the input is the
two-lattice comparison itself, which fails for $C_R\ne C_I$.)
\end{remark}

\subsection{Rotation invariance}

\begin{theorem}[invariance under the generators, Wick sector]\label{rot:thm:gen}
Let $a$ be $\mathrm W$-admissible ($\Wadm{a}$), and assume \Ka\ and \TIa. Then for every $R\in\Rgen$ and every $(f,g)\in\cK$,
\[
Z^a_\infty(\mathrm i\,R\cdot f,R\cdot g)=Z^a_\infty(\mathrm if,g),
\]
and consequently the same holds for every $R\in\Gz$.
\end{theorem}

\begin{proof}
Fix $R\in\Rgen$ and $(f,g)\in\cK$ supported in $\Qb\rho$.

\emph{Step 0 (geometry).} Take $L$ with $ML\ge\frac75\rho+2$. By \cref{rot:lem:geom}(c), $R\cdot(f,g)$ is supported in $\Qb{7\rho/5}$,
at $\ell^\infty$-distance $>1$ from $\R^3\setminus C_L$, and $(f,g)$ is supported in $\Qb{5ML/7}\subset R^{-1}\Qb{ML}$. Let
$\fc\in\Aw{R^{-1}\Qb{ML}}$. Then $R^{-1}\Lk$ descends to $D_{J'}$ (\cref{rot:lem:geom}(b)), and $B'_k\defeq R^{-1}(C_L\cap\Lk)$ is a set
of sites of $R^{-1}\Lk$ which, with its nearest neighbours, lies in the interior of the cell, and on which $w=0$. Since $R\cdot(f,g)$ has
distance $>1$ from $\R^3\setminus C_L$, every plaquette and site of $R^{-1}\Lk$ at which $(f,g)$ does not vanish lies in $B'_k$.

\emph{Step 1 (cutting on the rotated lattice).} By \cref{rot:lem:cutwall}, applied to the walled torus state
$\langle\cdot\rangle_{\fc,R^{-1}\Lk}$ with $X=B'_k$, whose $X$-factor is the Neumann--Dirichlet state on $B'_k$ with the scalar alive on
all of $B'_k$, and by the multiplier being the same constant on both sides (\cref{rot:lem:walledalg}(c), \cref{lem:s02:alg}(a)), at every $k$,
\[
Z^a_{R^{-1}\Lk,B'_k}(\mathrm if,g)\le\Tw{a}{\fc,R^{-1}\Lk}(f,g).
\]

\emph{Step 2 (passive rotation).} By \cref{rot:lem:walledalg}(b) with $O=R$, which maps $R^{-1}\Lk$ with its frame $(R^{-1}e_\alpha)$ onto
$\Lk$ with the frame $(e_\alpha)$ and $B'_k$ onto $X_L$,
$Z^a_{R^{-1}\Lk,B'_k}(\mathrm if,g)=Z^a_{k,L}(\mathrm i\,R\cdot f,R\cdot g)$.

\emph{Step 3 (continuum limits).} As $k\to\infty$ the left side converges by \cref{rot:thm:full}(i) (the sources $R\cdot(f,g)$ being at
distance $>1$ from $\R^3\setminus C_L$) and the right side by \Ka\ (\cref{rot:def:walls}):
\[
Z^a_L(\mathrm i\,R\cdot f,R\cdot g)\le\Wv^a_\fc(f,g)\qquad\text{for every }\fc\in\Aw{R^{-1}\Qb{ML}}.
\]

\emph{Step 4 (infimum, volume limit, closing hypothesis).} Take the infimum over $\fc\in\Aw{R^{-1}\Qb{ML}}$ and let $L\to\infty$. By
\cref{rot:thm:full}(iii), \cref{rot:lem:Rind} and \TIa,
\[
Z^a_\infty(\mathrm i\,R\cdot f,R\cdot g)=\lim_{L\to\infty}Z^a_L(\mathrm i\,R\cdot f,R\cdot g)\le\Wenva(f,g)\le Z^a_\infty(\mathrm if,g).
\]

\emph{Step 5 (the inverse).} $R^{-1}\in\Rgen$ (\cref{rot:lem:geom}(a)) and $R\cdot(f,g)\in\cK$, so Steps 0--4 for $(R^{-1},R\cdot(f,g))$
give $Z^a_\infty(\mathrm if,g)\le Z^a_\infty(\mathrm i\,R\cdot f,R\cdot g)$. Since $\cK$ is $\Gz$-invariant and the transport is a left
action, the claim extends to $\Gz$.
\end{proof}

\begin{remark}[the order of limits]\label{rot:rem:order}
The order of limits is: at fixed $R$, $(f,g)$, $L$ and $\fc$, first $k\to\infty$; then the infimum over $\fc$; then $L\to\infty$. The
continuum limit is always taken first, at fixed volume and fixed wall, as in King's construction~\cite[p.~336]{K2}, and the torus
comparison is used only as an equality of continuum limits, never at a fixed lattice spacing.
\end{remark}

\begin{theorem}[closure: from the generators to $O(3)$, Wick sector]\label{rot:thm:closure}
Let $a$ be $\mathrm W$-admissible ($\Wadm{a}$), and suppose $Z^a_\infty(\mathrm i\,R\cdot f,R\cdot g)=Z^a_\infty(\mathrm if,g)$ for all $(f,g)\in\cK$
and $R\in\Rgen$. Then for every $R\in O(3)$, every $N$, all real $2$-forms $f_j$ and real $h_j$ in $C^\infty_c$, and all
$(z,w)\in\C^{2N}$,
\[
Z^a_\infty\Big(\sum_jz_j\,R\cdot f_j,\sum_jw_j\,R\cdot h_j\Big)=Z^a_\infty\Big(\sum_jz_jf_j,\sum_jw_jh_j\Big).
\]
\end{theorem}

\begin{proof}
\emph{Complex coefficients.} The hypothesis passes to $\Gz$, since $\cK$ is invariant under $O(3)$ and $\Rgen=\Rgen^{-1}$. For
$R\in\Gz$ and a finite family, split $h_j=h^+_j-h^-_j$ as in the proof of \cref{rot:thm:full}(iii). The difference of the two sides, as a
function of $(z,w^+,w^-)\in\C^{3N}$, is entire by \cref{rot:thm:full}(iii) and vanishes on
$\Omega_+=(\mathrm i\R)^N\times(0,\infty)^{2N}$, where all sources lie in $\cK$; $\Omega_+$ has non-empty interior within a totally
real subspace of maximal dimension, hence a set of uniqueness \cite{Hor90}, and the difference vanishes identically.

\emph{Closure in $SO(3)$.} At real coefficients, $\Phi^a(R\cdot s)=\Phi^a(s)$ for $s\in\cE_\rho$ (the pairs of real sources supported
in $\Bbar\rho$, \cref{rot:thm:full}(iv)) and $R\in\Gz$; note $R\cdot s\in\cE_\rho$ since $R\Bbar\rho=\Bbar\rho$. The map
$R\mapsto R\cdot s\in\cE_\rho$ is continuous on $O(3)$, and $\Phi^a$ is continuous on $\cE_\rho$ by \cref{rot:thm:full}(iv). Hence
$G\defeq\{R\in O(3):\Phi^a(R\cdot s)=\Phi^a(s)\text{ for all }\rho\text{ and }s\in\cE_\rho\}$ is a closed subgroup of $O(3)$ containing
$\Gz$, so $G\supset SO(3)$ by \cref{rot:lem:geom}(d).

\emph{Reflection.} Let $\sigma_3=\diag(1,1,-1)$ and $(f,g)\in\cK$ supported at $\ell^\infty$-distance $>1$ from
$\R^3\setminus C_{J-1}$. \Cref{lem:s02:refl} gives, at every $k$,
$Z^a_{k,J-1}(\mathrm if,g)\le Z^a_{k,J}(\mathrm i\,\sigma_3\cdot f,\sigma_3\cdot g)\le Z^a_{k,J+1}(\mathrm if,g)$. The family
$\sigma_3\cdot(f,g)$ is supported at distance $>1$ from $\R^3\setminus\CJ$, the region $\{\dist_\infty(\cdot,\R^3\setminus C_{J-1})>1\}$
being $\sigma_3$-symmetric and contained in $\CJ$; so all three terms converge as $k\to\infty$ by \cref{rot:thm:full}(i), and then as
$J\to\infty$ by \cref{rot:thm:full}(iii): $Z^a_\infty(\mathrm i\,\sigma_3\cdot f,\sigma_3\cdot g)=Z^a_\infty(\mathrm if,g)$ on $\cK$. The
first step, with $\sigma_3$ in place of $R$, gives $\sigma_3\in G$, so $G=O(3)=SO(3)\cup\sigma_3SO(3)$.

\emph{Conclusion.} For $R\in O(3)$ both sides of the claim are entire in $(z,w)\in\C^{2N}$ (\cref{rot:thm:full}(iii)) and agree on
$\R^{2N}$, since $R\in G$; hence they agree on $\C^{2N}$.
\end{proof}

\begin{theorem}[the curvature sector]\label{rot:thm:curv}
King's printed state on curvature sources is
$O(3)$-invariant: for every $R\in O(3)$, all real $2$-forms $f_1,\dots,f_N\in C^\infty_c$ and all $z\in\C^N$,
\[
\ZKinf\Big(\sum_jz_j\,R\cdot f_j,0\Big)=\ZKinf\Big(\sum_jz_jf_j,0\Big),
\]
and the curvature Schwinger functions $S^K_n$ of \cref{thm:s05:CK}(c) are $O(3)$-covariant, $F$ transforming as a $2$-form. Moreover,
for every Euclidean motion $x\mapsto Rx+b$, $\ZKinf(\sum_jz_j\,\tau_bR\cdot f_j,0)=\ZKinf(\sum_jz_jf_j,0)$.
\end{theorem}

\begin{proof}
The inputs are the existence theorem \EC\ (\Cref{hyp:s03:EC}), the torus comparison \Kz\ (\Cref{rot:hyp:K}) and the closing estimate \TIz\ (\Cref{rot:hyp:TIGHT}), the last two on the class of \cref{rot:rem:class}.

\emph{Invariance under the generators at imaginary points.} Fix $R\in\Rgen$ and a real $2$-form $f\in C^\infty_c$ supported in
$\Qb\rho$. At $g=0$ the walled values are those of \Kz\ (\cref{rot:def:walls}). Steps 0--2 of the proof of \cref{rot:thm:gen} run
unchanged with the source $f$ alone: with $L$, $\fc\in\Aw{R^{-1}\Qb{ML}}$ and $B'_k$ as there, \cref{rot:lem:cutwall} on the rotated
lattice and \cref{rot:lem:walledalg}(b) give, at every $k$,
\[
Z_{k,L}(\mathrm i\,R\cdot f)=\ava{e^{\mathrm iF(R\cdot f)}}\big|_{J=L}=Z_{R^{-1}\Lk,B'_k}(\mathrm if)\le T_{\fc,R^{-1}\Lk}(f,0),
\]
the middle term being the functional of the comparison cube $X_L$. \emph{Step 3.} The right side converges to $\Wv_\fc(f,0)$ by \Kz\
(\cref{rot:lem:K}(a)), so $\limsup_kZ_{k,L}(\mathrm i\,R\cdot f)\le\Wv_\fc(f,0)$; and since $R\cdot f$ is supported at distance $>1$ from
$\R^3\setminus C_L$ (Step 0), \cref{cor:s04:G2}(ii)($\gamma$) under \EC\ gives $Z^K_L(\mathrm i\,R\cdot f)\le\liminf_kZ_{k,L}(\mathrm i\,R\cdot f)$.
Hence
\[
Z^K_L(\mathrm i\,R\cdot f)\le\Wv_\fc(f,0)\qquad\text{for every }\fc\in\Aw{R^{-1}\Qb{ML}}\text{ of the class}.
\]
\emph{Step 4.} Take the infimum over $\fc$ and let $L\to\infty$: the left side converges to $\ZKinf(\mathrm i\,R\cdot f,0)$ by
\cref{cor:s04:G2}(ii)($\beta$), the infima converge to $\Wenv(f,0)$ by \cref{rot:lem:Rind} (on the class, \cref{rot:rem:class}), and
\TIz\ gives
\[
\ZKinf(\mathrm i\,R\cdot f,0)\le\Wenv(f,0)\le\ZKinf(\mathrm if,0).
\]
\emph{Step 5.} $R^{-1}\in\Rgen$ and $R\cdot f$ is a real $2$-form in $C^\infty_c$, so the same for $(R^{-1},R\cdot f)$ gives the reverse
inequality: $\ZKinf(\mathrm i\,R\cdot u,0)=\ZKinf(\mathrm iu,0)$ for every real $2$-form $u\in C^\infty_c$ and $R\in\Rgen$, hence for
every $R\in\Gz$, the transport being a left action.

\emph{Closure in $SO(3)$.} Fix a real $2$-form $u\in C^\infty_c$. The map $R\mapsto R\cdot u$ is continuous from $O(3)$ to
$L^2(\R^3;\Lambda^2)$: for $R\to R_0$ the $2$-forms $R\cdot u$ converge to $R_0\cdot u$ uniformly, with supports in one compact set. By the
Lipschitz bound of \cref{cor:s06:TK}, $|\ZKinf(\mathrm iR\cdot u)-\ZKinf(\mathrm iR_0\cdot u)|\le\sqrt2\,\norm{R\cdot u-R_0\cdot u}_2$, so
$R\mapsto\ZKinf(\mathrm i\,R\cdot u,0)$ is continuous on $O(3)$. It is constant on $\Gz$, which is dense in $SO(3)$
(\cref{rot:lem:geom}(d)); hence $\ZKinf(\mathrm i\,R\cdot u,0)=\ZKinf(\mathrm iu,0)$ for all $R\in SO(3)$.

\emph{The reflection.} By \cref{thm:s06:S1K} with $O=\sigma_3\in\Oh$, $\ZKinf(\mathrm i\,\sigma_3\cdot u,0)=\ZKinf(\mathrm iu,0)$ for every
real $u$. Since $O(3)=SO(3)\cup\sigma_3SO(3)$ and $(\sigma_3R)\cdot u=\sigma_3\cdot(R\cdot u)$, the identity at imaginary points holds for
every $R\in O(3)$.

\emph{All complex coefficients.} Fix $R\in O(3)$ and the family $f_1,\dots,f_N$. By \cref{thm:s05:CK}(b) both sides of the claim are
entire functions of $z\in\C^N$. At $z=\mathrm it$, $t\in\R^N$, we have $\sum_j\mathrm it_j\,R\cdot f_j=\mathrm i\,R\cdot u_t$ with
$u_t=\sum_jt_jf_j$, so the two sides agree on $(\mathrm i\R)^N$, a totally real subspace of maximal dimension, hence a set of
uniqueness \cite{Hor90}; they agree on $\C^N$. Differentiating at $z=0$ gives $S^K_n(R\cdot f_1,\dots,R\cdot f_n)=S^K_n(f_1,\dots,f_n)$
for every $n$, the $O(3)$-covariance of the Schwinger functions.

\emph{Euclidean motions.} By \cref{cor:s06:TK}, $\ZKinf(\sum_jz_j\tau_bR\cdot f_j,0)=\ZKinf(\sum_jz_jR\cdot f_j,0)$ for every
$b\in\R^3$, and the right side equals $\ZKinf(\sum_jz_jf_j,0)$ by the above.
\end{proof}

\begin{corollary}[the curvature two-point function]\label{rot:cor:twopoint}
The two-point Schwinger function
$S^K_2(f_1,f_2)=\lim_J\lim_k\avK{F(f_1)F(f_2)}$ of \cref{thm:s05:CK}(c) satisfies $S^K_2(R\cdot f_1,R\cdot f_2)=S^K_2(f_1,f_2)$ for all
$R\in O(3)$ and all real $2$-forms $f_1,f_2\in C^\infty_c$.
\end{corollary}

\begin{proof}
Of the closing estimate only \TItwo\ is used. Let $R\in\Rgen$ and $f$ a real $2$-form. Steps 0--3 of the proof of \cref{rot:thm:curv}, with the source $tf$, give
$Z^K_L(\mathrm it\,R\cdot f)\le\Wv_\fc(tf,0)$ for all real $t$ and all $\fc\in\Aw{R^{-1}\Qb{ML}}$. Both sides are even and smooth in $t$
and equal $1$ at $t=0$; comparing second derivatives, $v_L(R\cdot f)\ge\sigma^2_\fc(f)$, where $v_L(u)\defeq\lim_k\avK{F(u)^2}|_{J=L}$
is the cube variance (\cref{thm:s05:CK}(c)). Since $\Aw{\Qb{7ML/5}}\subset\Aw{R^{-1}\Qb{ML}}$ and $b\mapsto\sup_{\Aw{\Qb b}}\sigma^2_\fc(f)$
is nonincreasing,
\[
v_L(R\cdot f)\ge\sup_{\Aw{\Qb{7ML/5}}}\sigma^2_\fc(f)\ge\lim_{b\to\infty}\sup_{\Aw{\Qb b}}\sigma^2_\fc(f)\ge S^K_2(f)
\]
by \TItwo. The cube variances converge to $S^K_2(R\cdot f)$ as $L\to\infty$ (\cref{thm:s05:CK}(c)), so $S^K_2(R\cdot f)\ge S^K_2(f)$, and
the inverse generator gives equality; hence $S^K_2(u,u)$ is $\Gz$-invariant.

\emph{The bound.} By \cref{thm:s05:CK}(a), $\avK{e^{tF(f)}}\le\exp(\frac12t^2\lmu^{-2}\norm{\dstar_kf}^2_{2,k})$ for real $t$, with equality
at $t=0$; comparing second derivatives, $\avK{F(f)^2}\le\lmu^{-2}\norm{\dstar_kf}^2_{2,k}$, and $\norm{\dstar_kf}_{2,k}\to\norm{\dstar f}_{L^2}$
by \eqref{eq:s02:riemann}. Letting $k\to\infty$ and then $J\to\infty$, $S^K_2(f)\le\lmu^{-2}\norm{\dstar f}^2_{L^2}$; and $S^K_2(f)\ge0$
since $\ZKinf(\mathrm itf,0)\le1=\ZKinf(0,0)$ (\cref{thm:s05:CK}(b)). (The continuum bound of \cref{thm:s05:CK}(b) gives the same at once.)

\emph{Closure.} $S^K_2$ is a positive semidefinite symmetric bilinear form, so
$|S^K_2(f,f)-S^K_2(f',f')|=|S^K_2(f-f',f+f')|\le S^K_2(f-f')^{1/2}S^K_2(f+f')^{1/2}$, and with the bound, $R\mapsto S^K_2(R\cdot f,R\cdot f)$
is continuous on $O(3)$, because $\dstar(R\cdot f)\to\dstar(R_0\cdot f)$ in $L^2$ as $R\to R_0$ (uniform convergence of the first
derivatives on one compact set). Density of $\Gz$ gives $SO(3)$; \cref{thm:s06:S1K} with $O=\sigma_3$, differentiated twice at $z=0$,
gives $\sigma_3$, hence $O(3)$; and polarisation gives the bilinear form.
\end{proof}
